\subsection{无限基数的性质}

定理2. 对每一个无限基数 \(\mathfrak{a}\) 我们拥有 \(\mathfrak{a}^2 = \mathfrak{a}\) .

我们将使用以下两个引理：

\textbf{引理1. 每个无限集E都包含一个与N等势的集合。}

在 E 上存在一个良序关系（§ 2，第 3 号，定理 1），我们将其记为 \(x \leqslant y\) 。该假设蕴含良序集 E 不能同构于 N 的一个不同于 N 的段，因为这样的段具有形式 {[}0, n{]} (§ 2, 第 1 号，命题 1) 因此是有限的。由此可知，N 同构于 E 的一个段（§ 2, 第 3 号，定理 3），从而得出结果。

\textbf{引理2. 集合N × N与N等势。}

由于 \(N \times N\) 包含集合 \(\{0\} \times N\) ，该集合与 N 等势，因此我们有 \(\text{Card}(N) \leqslant \text{Card}(N \times N)\) 。为了完成证明，只需定义从 \(N \times N\) 到 N 的一个单射 f。为此，我们注意到存在一个单射 \(\varphi\) （从 \(\mathbf{N}\) 到由 \(\mathbf{N}\) 到 \(\mathrm{I} = \{0,1\}\) 中的映射所组成的集合中），按如下方式获得：如果 \(r\) 是使得 \(n > 2^r\) 的最小整数，并且如果 \(\sum_{k=0}^{r-1}\varepsilon_k2^{r-k-1}\) 是 \(n\) 的二进展开 ( \(\S 5\) ，第7号）， \(\varphi(n)\) 被定义为序列 \((u_m)_{m\in\mathbf{N}}\) 使得 \(u_m = \varepsilon_{r-m-1}\) 对于 \(m < r\) 并且 \(u_m = 0\) 对于 \(m \geqslant r\) 。 \(\S 5\) 第7号的命题8表明 \(\varphi\) 是单射。对于每一对 \((n, n') \in \mathbf{N} \times \mathbf{N}\) ，我们定义 \(f(n, n')\) 如下：如果 \(\varphi(n) = (u_m)\) 并且 \(\varphi(n') = (v_m)\) ，让 \(f(n, n')\) 是整数 \(s\) 使得 \(\varphi(s) = (w_m)\) ，其中 \(w_{2m} = u_m\) 并且 \(w_{2m+1} = v_m\) 对于所有 \(m \in \mathbf{N}\) 。显然，这个关系 \(f(n, n') = f(n_1, n_1')\) 蕴含 \(\varphi(n) = \varphi(n_1)\) 并且 \(\varphi(n') = \varphi(n_1')\) ；因此 \((n, n') = (n_1, n_1')\) ，从而 \(f\) 是单射的。

现在我们来看定理2的证明。设E是一个集合，使得 \(\operatorname{Card}(\mathbf{E}) = \mathfrak{a}\) 。设D是E的一个与N等势的子集（引理1）。则存在一个双射 \(\psi_0\) 从D到D × D（引理2）。设 \(\mathfrak{M}\) 为所有对 (X, \(\psi\) ) ——其中 X 是 E 的一个子集，包含 D ，而 \(\psi\) 是X到X × X上扩展 \(\psi_0\) 的一个双射——所组成的集合. 通过关系对集合 \(\mathfrak{M}\) 进行排序

\textbf{“X ⊂ X' 且 ψ' 是 ψ 的扩张”}

在(X, \(\psi\) )与(X', \(\psi'\) )之间。则容易看出 \(\mathfrak{M}\) 是归纳的（参见§2，第4号，例2）。因此，由Zorn引理（§2，第4号，定理2）， \(\mathfrak{M}\) 有一个极大元(F, \(f\) )。我们将证明Card (F) = a，这足以证明定理。如果Card (F)不等于a，设Card (F) = b \textless{} a。则b = b²且b是无限的，且b ≤ 2b ≤ 3b ≤ b² = b（§3，第6号，命题14）；因此2b = b且3b = b。由假设b \textless{} a可知Card (E --- F) \textgreater{} b，因为否则我们会有Card (E) ≤ 2b = b，而我们已经假设b \textless{} Card (E)。因此存在一个子集Y ⊂ E --- F与F等势。设Z = F ∪ Y；我们将证明存在一个双射g从Z到Z × Z，它扩张了f。～我们有

\[
\mathbf {Z} \times \mathbf {Z} = (\mathbf {F} \times \mathbf {F}) \cup (\mathbf {F} \times \mathbf {Y}) \cup (\mathbf {Y} \times \mathbf {F}) \cup (\mathbf {Y} \times \mathbf {Y}),
\]

且右边的四个乘积两两不相交。由于F和Y等势，我们有

\[
\operatorname{Card} (\mathbf {F} \times \mathbf {Y}) = \operatorname{Card} (\mathbf {Y} \times \mathbf {F}) = \operatorname{Card} (\mathbf {Y} \times \mathbf {Y}) = \mathfrak {b} ^ {2} = \mathfrak {b},
\]

由此

\[
\operatorname{Card} ((\mathbf {F} \times \mathbf {Y}) \cup (\mathbf {Y} \times \mathbf {F}) \cup (\mathbf {Y} \times \mathbf {Y})) = 3 \mathfrak {b} = \mathfrak {b}.
\]

因此存在一个双射 \(f_{1}\) 从Y到集合

\[
(\mathbf {F} \times \mathbf {Y}) \cup (\mathbf {Y} \times \mathbf {F}) \cup (\mathbf {Y} \times \mathbf {Y});
\]

，映射g从Z到 \(Z \times Z\) ，它在F上等于f，在Y上等于 \(f_{1}\) ，因此是一个扩张f的双射，这与f的定义矛盾。～因此 \(\text{Card}(F) = a\) ，证明完成。

推论1。若 \(\mathfrak{a}\) 是一个无限基数，则 \(\mathfrak{a}^n = \mathfrak{a}\) 对于每一个整数 \(n \geqslant 1\) . 通过对 \(n\) 进行归纳。

推论2。一个有限族 \((\mathfrak{a}_i)_{i\in \mathbf{I}}\) 的非零基数的乘积，其中最大的是一个无限基数 \(\mathfrak{a}\) ，等于 \(\mathfrak{a}\) .

让 \(\mathfrak{b}\) 表示乘积，且设 \(n\) 为I中元素的个数。则 \(\mathfrak{b} \leqslant \mathfrak{a}^n = \mathfrak{a}\) （§3，第6号，命题14）。另一方面，由于 \(\mathfrak{a}_i \geqslant 1\) 对于所有 \(i \in \mathbf{I}\) ，我们有 \(\mathfrak{b} \geqslant \mathfrak{a}\) (§3，第6号，命题14)。

推论3。设 \(\mathfrak{a}\) 为一个无限基数，且设 \((\mathfrak{a}_{\iota})_{\iota \in I}\) 为一族基数 \(\leqslant \mathfrak{a}\) ，其指标集 \(I\) 具有基数 \(\leqslant \mathfrak{a}\) 。然后 \(\sum_{\iota \in I} \mathfrak{a}_{\iota} \leqslant \mathfrak{a}\) ；并且如果 \(\mathfrak{a}_{\iota} = \mathfrak{a}\) 对于至少一个指标 \(\iota \in I\) ，那么 \(\sum_{\iota \in I} \mathfrak{a}_{\iota} = \mathfrak{a}\) .

让 \(\mathfrak{b}\) 为 I 的基数；则我们有 \(\sum_{i\in I}a_i\leqslant a\mathfrak{b}\leqslant a^2 = a\) (§3，第6号，命题14)，以及 \(\sum_{i\in I}a_i\geqslant a_x\) 对于所有 \(x\in I\) .

推论4。如果a和b是两个非零基数，其中一个是无限的，我们有 \(a\mathfrak{b} = a + \mathfrak{b} = \sup(a, \mathfrak{b})\) .

这直接由推论2和推论3得出。

